% Part: first-order-logic
% Chapter: tableaux
% Section: propositional-rules

\documentclass[../../../include/open-logic-section]{subfiles}

\begin{document}

\iftag{FOL}
      {\olfileid{fol}{tab}{prl}}
      {\olfileid{pl}{tab}{prl}}

\olsection{विधानीय नियम}

\subsection{$\lnot$ साठीचे नियम}

\begin{defish}
\AxiomC{\sFmla{\True}{\lnot !A}}
\RightLabel{\TRule{\True}{\lnot}}
\UnaryInfC{\sFmla{\False}{!A}}
\DisplayProof
\hfill
\AxiomC{\sFmla{\False}{\lnot !A}}
\RightLabel{\TRule{\False}{\lnot}}
\UnaryInfC{\sFmla{\True}{!A}}
\DisplayProof
\end{defish}

\subsection{$\land$ साठीचे नियम}

\begin{defish}\noindent
\AxiomC{\sFmla{\True}{!A \land !B}}
\RightLabel{\TRule{\True}{\land}}
\UnaryInfC{\sFmla{\True}{!A}}
\noLine
\UnaryInfC{\sFmla{\True}{!B}}
\DisplayProof
\hfill
\AxiomC{\sFmla{\False}{!A \land !B}}
\RightLabel{\TRule{\False}{\land}}
\UnaryInfC{$\sFmla{\False}{!A} \quad \mid \quad \sFmla{\False}{!B}$}
\DisplayProof
\end{defish}

\subsection{$\lor$ साठीचे नियम}

\begin{defish}
\AxiomC{\sFmla{\True}{!A \lor !B}}
\RightLabel{\TRule{\True}{\lor}}
\UnaryInfC{$\sFmla{\True}{!A} \quad \mid \quad \sFmla{\True}{!B}$}
\DisplayProof
\hfill
\AxiomC{\sFmla{\False}{!A \lor !B}}
\RightLabel{\TRule{\False}{\lor}}
\UnaryInfC{\sFmla{\False}{!A}}
\noLine
\UnaryInfC{\sFmla{\False}{!B}}
\DisplayProof
\end{defish}

\subsection{$\lif$ साठीचे नियम}

\begin{defish}
\AxiomC{\sFmla{\True}{!A \lif !B}}
\RightLabel{\TRule{\True}{\lif}}
\UnaryInfC{$\sFmla{\False}{!A} \quad \mid \quad \sFmla{\True}{!B}$}
\DisplayProof
\hfill
\AxiomC{\sFmla{\False}{!A \lif !B}}
\RightLabel{\TRule{\False}{\lif}}
\UnaryInfC{\sFmla{\True}{!A}}
\noLine
\UnaryInfC{\sFmla{\False}{!B}}
\DisplayProof
\end{defish}

\subsection{कट नियम}

\begin{defish}
\AxiomC{}
\RightLabel{\Cut}
\UnaryInfC{$\sFmla{\True}{!A} \quad \mid \quad \sFmla{\False}{!A}$}
\DisplayProof
\end{defish}

\Cut{} नियम पूर्वीच्या !!{signed formula} वर ``लावला'' जात नाही; त्याऐवजी
तो !!a{tableau} मधील प्रत्येक शाखा दोन शाखांत फोडू देतो: एका शाखेत
$\sFmla{\True}{!A}$ आणि दुसरीत~$\sFmla{\False}{!A}$ असते. हा नियम आवश्यक
नाही---ज्या !!{signed formula}s च्या संचासाठी बंद !!{tableau} आहे, त्या
संचासाठी \Cut{} न वापरणाराही एक बंद टॅब्लो असतो---परंतु त्यामुळे
!!{tableau}s सोयीस्कर रीतीने एकत्र जोडता येतात.

\end{document}
